\chapter*{Samenvatting}
Het \emph{kwantumveeldeeltjesprobleem} bevindt zich in het middelpunt van de theoretische fysica sinds het begin van de vorige eeuw. Het centrale probleem is het voorspellen en begrijpen van het \emph{fasediagram} van microscopische \emph{kwantumhamiltonianen} die een macroscopisch aantal interagerende vrijheidsgraden beschrijven. Omwille van de \emph{exponenti\"ele} schaling van de dimensie van de veeldeeltjes \emph{Hilbertruimte} is het typisch onmogelijk om analytische oplossingen te construeren voor meer dan een handvol deeltjes. Vooruitgang hierin vraagt daarom om \emph{effici\"ente} en \emph{schaalbare benaderende} parameterisaties van de relevante golffuncties.

\emph{Tensor netwerken} zijn voortgekomen uit het domein van de \emph{kwantuminformatietheorie}, en in het bijzonder uit de \emph{kwantum verstrengelingstheorie}. Ze hebben de afgelopen twintig jaar veel aandacht gekregen in de studie van lokale kwantumhamiltonianen in lage spatiale dimensies. Tensor netwerken parametriseren de klasse van golffuncties die voldoen aan een \emph{oppervlaktewetschaling} van de \emph{verstrengelingsentropie}, die wordt waargenomen dat ze geldt voor lokale Hamiltonianen met een \emph{energie kloof}. Het centrale idee is dat \emph{globale} eigenschappen van de golffunctie effici\"ent kunnen worden gecomprimeerd in \emph{lokale} tensoren die de variationele parameters van de golffunctie bevatten. Een van de belangrijkste voordelen van deze lokale parametrisatie is dat de aanwezigheid van \emph{symmetrie\"en} zich vertaalt in lokale invariantie-eigenschappen van de onderliggende tensoren.

\emph{Symmetrie}, en in het bijzonder \emph{symmetriebreking}, wordt sinds lang geapprecieerd als een organiserend principe voor klassieke fasetransities, vooraleer zij even centraal kwam te staan in de studie van kwantumveeldeeltjesfasediagrammen. Traditioneel werken symmetrie\"en in op de onderliggende kwantumvrijheidsgraden via unitaire of anti-unitaire representaties van \emph{groepen}. Binnen \emph{Landau's paradigma} van fasetransities wordt een kwantum fase van materie gekarakteriseerd door de \emph{ongebroken} symmetrie van de grondtoestandsdeelruimte.

In recente jaren werd echter erkend dat er fasen bestaan die voorbij het traditionele Landau paradigma gaan. Deze fasen worden gekarakteriseerd door het bestaan van \emph{verborgen} symmetrieoperatoren die inwerken op de verstrengelingspatronen van de grondtoestand. Omdat tensor netwerken directe toegang verschaffen tot deze verstrengelingsvrijheidsgraden, onthullen ze deze verborgen symmetrie\"en en vormen ze een natuurlijke taal voor hun studie. Opmerkelijk genoeg hoeven deze verborgen symmetrieoperatoren niet inverteerbaar te zijn en kunnen ze daarom buiten het kader van groepentheorie vallen. Ze worden daarentegen op natuurlijke wijze beschreven door \emph{(hogere) fusiecategorie\"en}. In zekere zin vormen tensor netwerken concrete \emph{representaties} van deze gegeneraliseerde symmetriestructuren.

Deze categorische symmetrie\"en blijken ook te worden overge\"erfd door de \emph{randtheorie\"en} van deze fasen voorbij het Landau paradigma. Deze randtheorie\"en leven in een dimensie lager. Deze \emph{holografische} manier om gegeneraliseerde noties van symmetrie te beschrijven, heeft de voorbije jaren aanzienlijke aandacht gekregen en vormt een van de centrale onderwerpen van deze thesis. Concreet laat dit moderne perspectief op globale symmetrie\"en ons toe om gevestigde technieken uit de \emph{topologische kwantumveldentheorie} te gebruiken in hun studie. Dit gezichtspunt faciliteert in het bijzonder de studie van het \emph{ijken} en \emph{dualiseren} van deze symmetrie\"en, maar ook de classificatie van de fasen die erdoor worden beschermd, en past deze fasen zo in een \emph{gegeneraliseerd Landau paradigma}.

\bigskip\noindent Deze thesis behandelt verschillende aspecten van gegeneraliseerde symmetrie, in het bijzonder de bijbehorende dualiteiten en anomalie\"en, binnen de context van tensor netwerken. Ze is verdeeld in twee delen, \hyperref[part:Foundations]{\emph{Foundations}} and \hyperref[part:Applications]{\emph{Applications}}, die respectievelijk inleidende hoofdstukken over de onderwerpen van deze thesis en nieuwe onderzoeksresultaten in de vorm van vijf publicaties en een verzameling cursusnota's bevatten.

In het inleidende hoofdstuk~\ref{chap:introduction} beginnen we met een gedetailleerde beschrijving van de uitdaging die het kwantumveeldeeltjesprobleem stelt. We recapituleren enkele van de belangrijkste eigenschappen van het traditionele Landau paradigma en de rol die symmetrie daarin speelt.

Hoofdstuk~\ref{ch:TN} is gewijd aan de introductie van tensor netwerktoestanden en enkele van hun eigenschappen. Zowel de eendimensionale matrixproducttoestanden als de tweedimensionale geprojecteerde verstrengelde-paartoestanden worden besproken. Door gebruik te maken van de fundamentele stelling van matrixproducttoestanden, die in het bijzonder beschrijft hoe een symmetrie van een MPS zich vertaalt in een lokale invariantie van de onderliggende tensoren, herhalen we in detail de classificatie van eendimensionale kwantumfasen van materie. Deze fasen worden volledig beschreven door het patroon van symmetriebreking in de grondtoestandsruimte, samen met een topologische invariant voor de overblijvende symmetrie, die de symmetrie-beschermde topologische orde karakteriseert.

In hoofdstuk~\ref{ch:topo} geven we eerst een introductie tot de nodige begrippen uit fusiecategorietheorie. Vervolgens herhalen we de constructie van de string-net modellen van Levin en Wen. Deze vormen een kader voor roosterrealisaties van niet-chirale topologische ordes in (2+1) dimensies met randvoorwaarden die de energie kloof behouden. We bespreken in het bijzonder de classificatie van deze randvoorwaarden, evenals hun grondtoestanden en het anyonische excitatiespectrum. Cruciaal is dat deze grondtoestanden exacte tensor netwerkrepresentaties hebben die worden gekarakteriseerd door topologische symmetrie\"en die uitsluitend inwerken op de verstrengelingsvrijheidsgraden. Deze vormen precies de eerder vermelde verborgen symmetrie\"en.

Tot slot construeren we in hoofdstuk~\ref{ch:GenSym} algemene randhamiltonianen voor de string-net modellen. Deze Hamiltonianen erven de (in het algemeen niet-inverteerbare) symmetrie\"en van de bulk over. Vervolgens schetsen we hoe deze symmetrie\"en kunnen worden geijkt of gedualiseerd, en illustreren we dit aan de hand van de Kramers-Wannier-dualiteit van het Ising model in een transversaal veld. Op periodieke randvoorwaarden gedragen dualiteitstransformaties zich als niet-inverteerbare lineaire operatoren. Ze kunnen worden gepromoveerd tot isometrie\"en door de Hilbertruimte uit te breiden met bijkomende vrijheidsgraden die symmetrie-compatibele randvoorwaarden implementeren, waarop de dualiteitsoperatoren niet-triviaal kunnen inwerken. Vervolgens bestuderen we de werking van de tweedimensionale Kramers-Wannier-dualiteit en beargumenteren we dat deze leidt tot duale modellen met een symmetriestructuur die op natuurlijke wijze wordt beschreven door een fusie-2-categorie. Na een schets van het relevante algebra\"ische kader besluiten we dit hoofdstuk met een vooruitblik op een algemene theorie van dualiteiten in (2+1) dimensies.

Het eerste hoofdstuk van het tweede deel bestaat uit cursusnota's over tensor netwerken. Dit is het enige toepassingshoofdstuk dat geen fundamenteel nieuwe onderzoeksresultaten bevat. Bovendien zijn delen van de inleidende hoofdstukken gebaseerd op relevant materiaal uit deze cursusnota's.

De eerste publicatie, gepresenteerd in hoofdstuk~\ref{app:frieze}, classificeert eendimensionale bosonische fasen van materie die worden beschermd door quasi-eendimensionale ruimtelijke symmetrie\"en. Deze symmetrie\"en corresponderen met de zogenaamde frieze symmetriegroepen. Het hoofdstuk gebruikt technieken die sterk gelijken op die uit hoofdstuk~\ref{ch:TN} en levert, naast andere resultaten, canonieke vormen voor de fases die uit deze classificatie voortkomen. We geven ook een praktisch algoritme voor het berekenen van groepencohomologie en representatieve cocycles.

In de publicatie die in hoofdstuk~\ref{app:SN} wordt gepresenteerd, tonen we aan dat verschillende Morita-equivalente string-net modellen dezelfde niet-chirale topologische fase realiseren. Hoewel deze string-net modellen microscopisch sterk verschillen, construeren we een expliciet kwantumcircuit van constante diepte dat het ene model afbeeldt op het andere. Dit circuit werkt niet alleen in op de grondtoestandsdeelruimte, maar ook op ge\"exciteerde toestanden wanneer bijkomende hulpvrijheidsgraden in rekening worden gebracht.

Hoofdstuk~\ref{app:gauging} toont aan dat de eendimensionale dualiteitstransformaties die in hoofdstuk~\ref{ch:GenSym} worden ge\"introduceerd, kunnen worden ge\"interpreteerd als het ijken van een specifieke collectie symmetrie\"en uit een symmetriefusiecategorie. In het bijzonder tonen we aan dat de MPO dualiteitsoperatoren worden verkregen uit een ijkingsafbeelding waarop vervolgens een unitair kwantumcircuit van constante diepte wordt toegepast dat de ijkvrijheidsgraden ontstrengelt. Dit veralgemeent een eerdere constructie van Haegeman et al. voor groepssymmetrie\"en. We bespreken ook de rol van anomalie\"en binnen dit kader en bepalen onder welke voorwaarden deze ijkingsprocedure kan worden beschouwd als de implementatie van een lokale versie van de symmetrie.


In hoofdstuk~\ref{app:codes} bouwen we voort op een eerdere constructie van Jos\'e Garre-Rubio. We vertrekken van een abelse symmetrie in een arbitraire dimensie en ijken op iteratieve wijze zowel de symmetrie als haar duale symmetrie. Zo bekomen we toestanden in een dimensie hoger. We tonen aan dat de aldus verkregen toestanden op natuurlijke wijze stabilisator codes defini\"eren en geven verschillende voorbeelden die binnen dit kader passen.

Tot slot gaan we in hoofdstuk~\ref{app:TwistedGauging} over tot dualiteiten in (2+1) dimensies. We bespreken een dualiteitstransformatie in zowel een als twee ruimtelijke dimensies die neerkomt op de compositie van een ijkingstransformatie en een daaropvolgende getwiste ijkingstransformatie. Dit project lost het open probleem op van de constructie van zelfduale modellen in (2+1) dimensies, waarmee we modellen bedoelen die invariant zijn onder een specifieke dualiteitstransformatie. We promoveren deze dualiteitstransformaties bovendien tot unitaire transformaties door de gepaste randvoorwaarden in rekening te brengen.